\documentclass[10pt,a4paper]{amsart}
\usepackage{amsmath,amssymb,amsthm,graphicx}
\usepackage{verbatim}
\newenvironment{reference}{\comment}{\endcomment}
%\newenvironment{reference}{}{}
\newenvironment{slogan}{\comment}{\endcomment}
\newenvironment{history}{\comment}{\endcomment}
% For commutative diagrams we use Xy-pic
\usepackage[all]{xy}

% We use 2cell for 2-commutative diagrams.
\xyoption{2cell}
% Theorem environments.
%
\theoremstyle{plain}
\newtheorem{theorem}[subsection]{Theorem}
\newtheorem{proposition}[subsection]{Proposition}
\newtheorem{lemma}[subsection]{Lemma}

\theoremstyle{definition}
\newtheorem{definition}[subsection]{Definition}
\newtheorem{example}[subsection]{Example}
\newtheorem{exercise}[subsection]{Exercise}
\newtheorem{situation}[subsection]{Situation}

\theoremstyle{remark}
\newtheorem{remark}[subsection]{Remark}
\newtheorem{remarks}[subsection]{Remarks}

\numberwithin{equation}{subsection}

% Macros
%
\def\lim{\mathop{\mathrm{lim}}\nolimits}
\def\colim{\mathop{\mathrm{colim}}\nolimits}
\def\Spec{\mathop{\mathrm{Spec}}}
\def\Hom{\mathop{\mathrm{Hom}}\nolimits}
\def\Ext{\mathop{\mathrm{Ext}}\nolimits}
\def\SheafHom{\mathop{\mathcal{H}\!\mathit{om}}\nolimits}
\def\SheafExt{\mathop{\mathcal{E}\!\mathit{xt}}\nolimits}
\def\Sch{\mathit{Sch}}
\def\Mor{\mathop{\mathrm{Mor}}\nolimits}
\def\Ob{\mathop{\mathrm{Ob}}\nolimits}
\def\Sh{\mathop{\mathit{Sh}}\nolimits}
\def\NL{\mathop{N\!L}\nolimits}
\def\CH{\mathop{\mathrm{CH}}\nolimits}
\def\proetale{{pro\text{-}\acute{e}tale}}
\def\etale{{\acute{e}tale}}
\def\QCoh{\mathit{QCoh}}
\def\Ker{\mathop{\mathrm{Ker}}}
\def\Im{\mathop{\mathrm{Im}}}
\def\Coker{\mathop{\mathrm{Coker}}}
\def\Coim{\mathop{\mathrm{Coim}}}

% Boxtimes
%
\DeclareMathSymbol{\boxtimes}{\mathbin}{AMSa}{"02}

%
% Macros for moduli stacks/spaces
%
\def\QCohstack{\mathcal{QC}\!\mathit{oh}}
\def\Cohstack{\mathcal{C}\!\mathit{oh}}
\def\Spacesstack{\mathcal{S}\!\mathit{paces}}
\def\Quotfunctor{\mathrm{Quot}}
\def\Hilbfunctor{\mathrm{Hilb}}
\def\Curvesstack{\mathcal{C}\!\mathit{urves}}
\def\Polarizedstack{\mathcal{P}\!\mathit{olarized}}
\def\Complexesstack{\mathcal{C}\!\mathit{omplexes}}
% \Pic is the operator that assigns to X its picard group, usage \Pic(X)
% \Picardstack_{X/B} denotes the Picard stack of X over B
% \Picardfunctor_{X/B} denotes the Picard functor of X over B
\def\Pic{\mathop{\mathrm{Pic}}\nolimits}
\def\Picardstack{\mathcal{P}\!\mathit{ic}}
\def\Picardfunctor{\mathrm{Pic}}
\def\Deformationcategory{\mathcal{D}\!\mathit{ef}}
\usepackage{cleveref}
\crefname{originallemma}{original lemma}{original lemmas}
\Crefname{originallemma}{Original lemma}{Original lemmas}
\makeatletter
\newlabel{lemma-internal-hom-exact}{{lemma-internal-hom-exact}{source}{}{sourcepointer-1}{}}
\newlabel{lemma-internal-hom-exact@cref}{{[originallemma][0][]lemma-internal-hom-exact}{[1][0][]source}{}{}{}}
\makeatother
\crefname{originalreference}{original reference}{original references}
\Crefname{originalreference}{Original reference}{Original references}
\makeatletter
\newlabel{sheaves-section-limits-sheaves}{{sheaves-section-limits-sheaves}{source}{}{sourcepointer-2}{}}
\newlabel{sheaves-section-limits-sheaves@cref}{{[originalreference][0][]sheaves-section-limits-sheaves}{[1][0][]source}{}{}{}}
\makeatother
\title{Homomorphisms out of a finitely presented module commute with filtered colimits}
\author{The Stacks Project Authors}
\date{stacks-project:modules.tex:local-snapshot-20261012; GNU FDL 1.2 or later}
\begin{document}
\maketitle
\noindent\textit{Editorial note.} All delivered mathematical text is the author's own, copied from the pinned Stacks Project chapter source sources/stacks-modules/pinned-original-source.tex; no statement or proof step was rewritten, reordered or removed.\par\medskip
\noindent\textit{Editorial note.} Preamble provenance: the Stacks Project keeps one single shared preamble for every chapter; this card takes it directly from the staged sources/stacks-modules/preamble.tex as three verbatim spans (lines 8-12, 14-18 and 154-217), so no derived file was written and nothing was added or redefined. The wrapper class and the base packages are supplied by the assembler.\par\medskip
\noindent\textit{Editorial note.} Bibliography provenance: neither the counted statement nor its proof contains a citation (this lane rescans the whole span before writing the card), so the card carries no bibliography and no entry was reconstructed.\par\medskip
\noindent\textit{Editorial note.} Selection provenance: this card was opened under the file-based control contract (CONTROL.md), with the mechanical pre-card gate run on its statement and proof interval before the card was written.\par\medskip
\noindent\textit{Allowed standard prerequisites (editorial).}\begin{itemize}
\item The compatibility of filtered colimits of modules with restriction to open subsets. Invoked at the start of the counted proof (source lines 3990-3991) as a cross-chapter pointer. Reduces the statement to the case of a module admitting a global presentation.
\item The exactness of the sheaf of homomorphisms in the first variable for a right exact sequence. Invoked inside the counted proof at source line 4006 as a pointer to a lemma of this chapter. Produces the two exact rows of the diagram whose outer vertical arrows are computed directly.
\end{itemize}\medskip
\ifdefined\hypertarget\hypertarget{sourcepointer-1}{}\fi
\noindent\textit{Source reference.} Original lemma lemma-internal-hom-exact: sources/stacks-modules/pinned-original-source.tex, source lines 3887--3887. Navigation pointer only: the counted text refers to the Stacks item lemma-internal-hom-exact, which belongs to another chapter of the project and is not part of this card. The pointer carries the target's own label as its number because this chapter's staged snapshot does not contain the other chapter's tag table; its content is not reproduced here.\par\medskip
\ifdefined\hypertarget\hypertarget{sourcepointer-2}{}\fi
\noindent\textit{Source reference.} Original reference sheaves-section-limits-sheaves: sources/stacks-modules/pinned-original-source.tex, source lines 2640--2640. Navigation pointer only: the counted text refers to the Stacks item sheaves-section-limits-sheaves, which belongs to another chapter of the project and is not part of this card. The pointer carries the target's own label as its number because this chapter's staged snapshot does not contain the other chapter's tag table; its content is not reproduced here.\par\medskip
\begin{lemma}
\label{lemma-hom-finite-presentation-colimit}
Let $X$ be a ringed space. Let $\mathcal{F}$ be an $\mathcal{O}_X$-module
of finite presentation. Let
$\mathcal{G} = \colim_{\lambda \in \Lambda} \mathcal{G}_\lambda$
be a filtered colimit of $\mathcal{O}_X$-modules. Then the canonical map
$$
\colim_\lambda \SheafHom_{\mathcal{O}_X}(\mathcal{F}, \mathcal{G}_\lambda)
\longrightarrow
\SheafHom_{\mathcal{O}_X}(\mathcal{F}, \mathcal{G})
$$
is an isomorphism.
\end{lemma}

\begin{proof}
Taking colimits of sheaves of modules commutes with restriction to opens, see
Sheaves, Section \ref{sheaves-section-limits-sheaves}. Hence we may assume
$\mathcal{F}$ has a global presentation
$$
\bigoplus\nolimits_{j = 1, \ldots, m}
\mathcal{O}_X
\longrightarrow
\bigoplus\nolimits_{i = 1, \ldots, n}
\mathcal{O}_X
\to
\mathcal{F}
\to
0
$$
The functor $\SheafHom_{\mathcal{O}_X}(-, -)$ commutes with finite direct sums
in either variable and $\SheafHom_{\mathcal{O}_X}(\mathcal{O}_X, -)$
is the identity functor. By this and by Lemma \ref{lemma-internal-hom-exact}
we obtain an exact sequence
$$
0 \to
\SheafHom_{\mathcal{O}_X}(\mathcal{F}, \mathcal{G}) \to
\bigoplus\nolimits_{i = 1, \ldots, n} \mathcal{G} \to
\bigoplus\nolimits_{j = 1, \ldots, m} \mathcal{G}
$$
Since filtered colimits are exact in $\textit{Mod}(\mathcal{O}_X)$
also the top row in the following commutative diagram is exact
$$
\xymatrix{
0 \ar[r] &
\colim_\lambda \SheafHom_{\mathcal{O}_X}(\mathcal{F}, \mathcal{G}_\lambda)
\ar[r] \ar[d] &
\colim_\lambda \bigoplus\nolimits_{i = 1, \ldots, n} \mathcal{G}_\lambda
\ar[r] \ar[d] &
\colim_\lambda \bigoplus\nolimits_{j = 1, \ldots, m} \mathcal{G}_\lambda
\ar[d] \\
0 \ar[r] &
\SheafHom_{\mathcal{O}_X}(\mathcal{F}, \mathcal{G}) \ar[r] &
\bigoplus\nolimits_{i = 1, \ldots, n} \mathcal{G} \ar[r] &
\bigoplus\nolimits_{j = 1, \ldots, m} \mathcal{G}
}
$$
Since the right two vertical arrows are isomorphisms we conclude.
\end{proof}

\end{document}
